\documentclass[11pt]{amsart}

\usepackage[T1]{fontenc}
\usepackage{lmodern}
\usepackage{microtype}
\usepackage{amsmath,amssymb,amsthm}
\usepackage[hidelinks]{hyperref}

\hypersetup{
  pdftitle={Mesh-pattern pairs 113-116 are not equidistributed with
  the X\textasciicircum(1)/Y\textasciicircum(1) family: a
  counterexample to an open problem of Lv and Zhang}
}

\newtheorem{theorem}{Theorem}
\theoremstyle{plain}
\newtheorem{conjecture}[theorem]{Conjecture}
\theoremstyle{remark}
\newtheorem{remark}[theorem]{Remark}

\newcommand{\Sn}{\mathfrak{S}_n}
\newcommand{\occ}{\operatorname{occ}}

\title[Mesh-pattern pairs $113$--$116$ are not equidistributed]
{Mesh-pattern pairs $113$--$116$ are not equidistributed with
the $X^{(1)}/Y^{(1)}$ family: a counterexample to an open problem of
Lv and Zhang}
\date{}

\subjclass[2020]{05A05, 05A15}
\keywords{mesh pattern, joint distribution, permutation pattern, counterexample}

\begin{document}

\begin{abstract}
In the concluding remark of their paper on joint equidistributions of
the mesh patterns $123$ and $132$ with minus antipodal shadings, Lv and
Zhang record an unnumbered assertion, suggested by computer experiments
and left there explicitly as an open problem: the four mesh-pattern
pairs labelled $113$--$116$ have the same bivariate occurrence
distribution as the pairs $X_i^{(1)}$ and $Y_i^{(1)}$, $9\le i\le18$.
We disprove it on $\mathfrak S_4$. The former pairs have no
permutation with occurrence vector $(2,0)$, whereas every pair in
the latter family has exactly one. More generally, for every
$n\ge4$ the respective counts are $0$ and at least $(n-1)!/6$. The
failure is immediate from Lv and Zhang's own Theorems~6.1 and~6.2,
in Section~6 of the same paper, so the assertion is best read as a
slip in the concluding remark rather than as a conjecture.
\end{abstract}

\maketitle

All pair labels follow Lv and Zhang~\cite{LvZhang}. For an ordered
pair $P=(p_1,p_2)$ of mesh patterns, set
\[
T^P_{n,k,\ell}
=
\bigl|\{\pi\in\Sn:
\occ_{p_1}(\pi)=k,\ \occ_{p_2}(\pi)=\ell\}\bigr|.
\]
Write $c(r,s)$ for the unsigned Stirling number of the first kind.
Let $\mathcal P$ be the four pairs labelled $113,114,115,116$, and let
\[
\mathcal Q
=
\{X_i^{(1)}:9\le i\le18\}
\cup
\{Y_i^{(1)}:9\le i\le18\}.
\]
Section~7 of~\cite{LvZhang}, headed ``Concluding remark'', contains the
following unnumbered sentence, in which equations~(1) and~(2) are the
two displays of Theorem~2.4 of that paper, giving respectively the
joint distribution of the pairs in $\mathcal Q$ and its generating
function: ``In fact, computer experiments suggest that pairs $113$ to
$116$ have the same joint distribution as pairs
$\{X^{(1)}_i\}_{i=9}^{18}$ and $\{Y^{(1)}_i\}_{i=9}^{18}$, satisfying
(1) and (2); we leave justifying this observation as an open problem.''
Formalised, and weakened to the part of it that we shall use, this is
the following statement. It is unnumbered in~\cite{LvZhang} and is not
the only cross-pair statement made in that section, so we do not give
it a name beyond its location.

\begin{conjecture}[{\cite[Section~7]{LvZhang}}, unnumbered]
\label{conj:cross}
For every $P\in\mathcal P$, $Q\in\mathcal Q$, $n\ge1$, and
$k,\ell\ge0$,
\[
T^P_{n,k,\ell}=T^Q_{n,k,\ell}.
\]
\end{conjecture}

Two remarks on the standing of this assertion. First, its failure is
immediate from Theorems~6.1 and~6.2 of~\cite{LvZhang}, proved in
Section~6 of that paper, which show that each component of each of the
pairs $113$--$116$ has at most one occurrence in any permutation: the
joint distribution of such a pair is supported on $\{0,1\}^2$, while
the distribution asserted for it, equation~(1) of~\cite{LvZhang}, is
not. Second, Conjecture~4 of~\cite{LvZhang}, stated in the same
section, asserts of the eight pairs $119$--$126$ exactly what the
sentence above asserts of $113$--$116$, namely that they ``have the
same joint distribution as $\{X^{(1)}_i\}_{i=9}^{18}$ and
$\{Y^{(1)}_i\}_{i=9}^{18}$''. The sentence is therefore best read as a
slip in the concluding remark rather than as a conjecture. What follows
records the refutation and the minimal order at which it takes effect.

\begin{theorem}
For every $P\in\mathcal P$, $Q\in\mathcal Q$, and $n\ge4$,
\[
T^P_{n,2,0}=0,
\qquad
T^Q_{n,2,0}
=
(n-1)!\sum_{j=2}^{n-1}\frac{c(j-1,2)}{j!}
\ge \frac{(n-1)!}{6}>0.
\]
At $n=4$, one has $T^Q_{4,2,0}=1$. Consequently,
Conjecture~\ref{conj:cross} is false, and $n=4$ is the first
order at which it fails.
\end{theorem}

\begin{proof}
The proofs of Theorems~6.1 and~6.2 of~\cite{LvZhang} show that each
component of the pairs labelled $113$ and $115$ occurs at most once
in any permutation. The same conclusion for $114$ and $116$ follows
by inversion. Hence
\[
T^P_{n,2,0}=0
\qquad(P\in\mathcal P).
\]

For every $Q\in\mathcal Q$, Theorem~2.4, equation~(1),
of~\cite{LvZhang} gives, whenever $(k,\ell)\ne(0,0)$,
\[
T^Q_{n,k,\ell}
=
(n-1)!\binom{k+\ell}{k}
\sum_{j=2}^{n-1}\frac{c(j-1,k+\ell)}{j!}.
\]
Therefore
\[
T^Q_{n,2,0}
=
(n-1)!\sum_{j=2}^{n-1}\frac{c(j-1,2)}{j!}.
\]
For $n\ge4$, the term $j=3$ is present, all terms are nonnegative,
and $c(2,2)=1$. Thus
\[
T^Q_{n,2,0}
\ge
\frac{(n-1)!}{3!}.
\]
For $n=4$, the only nonzero term is $j=3$, since $c(1,2)=0$; so
$T^Q_{4,2,0}=1$.

For $n\le2$, no length-three pattern occurs. For $n=3$, the only
three-term subsequence is the whole permutation, so the shading is
vacuous; every pair considered here has one permutation with vector
$(1,0)$, one with vector $(0,1)$, and four with vector $(0,0)$.
Hence $n=4$ is minimal.
\end{proof}

\begin{remark}[An explicit counterexample]
In the standard notation $(\sigma,R)$ for a mesh pattern, the pair
labelled $113$ and the pair $X_9^{(1)}$ are
\[
P_{113}
=
\bigl((123,R_{113}),(132,R_{113})\bigr),
\qquad
X_9^{(1)}
=
\bigl((123,R_9),(132,R_9)\bigr),
\]
where
\begin{align*}
R_{113}
&=
\{(0,1),(0,2),(1,1),(1,2),(2,2),
  (3,0),(3,1),(3,2),(3,3)\},\\
R_9
&=
\{(0,0),(0,1),(0,2),(0,3),
  (1,1),(1,2),(3,1),(3,2)\}.
\end{align*}
For $P_{113}$, the shaded rightmost column forces every occurrence
to end at the final entry. The shaded cells $(a,2)$,
$0\le a\le3$, then force the middle-position entry to equal the
final entry minus $1$ for the $123$ component, or plus $1$ for the
$132$ component. The first entry must be the largest entry to the
left of the middle-position entry that is smaller than it, and is
therefore unique. Thus each component occurs at most once.

For $X_9^{(1)}$, the permutation $1234$ has two occurrences of its
$123$ component, on positions $(1,2,3)$ and $(1,2,4)$. The unused
point lies in the unshaded cell $(3,3)$ for the first occurrence and
in the unshaded cell $(2,2)$ for the second; the other two increasing
triples are excluded by the shaded cells $(1,1)$ and $(0,0)$. Since
$1234$ has no classical $132$ occurrence, its occurrence vector is
$(2,0)$.

The discrepancy is not confined to the cell $(2,0)$. On
$\mathfrak S_4$ the joint distribution of $P_{113}$ is
\[
(0,0)\mapsto12,\;
(0,1)\mapsto5,\;
(1,0)\mapsto5,\;
(1,1)\mapsto2,
\]
while that of $X_9^{(1)}$ is
\[
(0,0)\mapsto12,\;
(0,1)\mapsto4,\;
(0,2)\mapsto1,\;
(1,0)\mapsto4,\;
(1,1)\mapsto2,\;
(2,0)\mapsto1,
\]
so the two already differ in the cells $(0,1)$ and $(1,0)$ as well.
The cell $(0,0)$, on the other hand, agrees: direct enumeration gives
$T^{P_{113}}_{n,0,0}=2(n-1)!$ for $2\le n\le7$, which is exactly the
value of the $k=\ell=0$ branch of equation~(1) of~\cite{LvZhang}. A
computation testing only avoidance would therefore not have detected
the failure.
\end{remark}

\begin{remark}[Scope]
The assertion refuted here is the unnumbered cross-pair statement in
Section~7 of~\cite{LvZhang}. It is distinct from the seven within-pair
symmetry conjectures involving fourteen labelled patterns in
Table~11, which were proved by Fang, Fu, Kitaev, and
Li~\cite{FangFuKitaevLi}. The other comparisons between different
pairs proposed by Lv and Zhang, namely Conjectures~3 and~4
of~\cite{LvZhang} on the pairs $117$, $118$ and $119$--$126$, are not
considered here.
\end{remark}

\begin{thebibliography}{9}

\bibitem{LvZhang}
S.~Lv and P.~B. Zhang,
\newblock Joint equidistributions of mesh patterns 123 and 132 with
minus antipodal shadings,
\newblock \emph{Discrete Appl. Math.} \textbf{379} (2026), 419--433,
\newblock
\href{https://doi.org/10.1016/j.dam.2025.09.002}
{doi:10.1016/j.dam.2025.09.002}.

\bibitem{FangFuKitaevLi}
Q.~Fang, S.~Fu, S.~Kitaev, and H.~Li,
\newblock On fourteen equidistribution conjectures of Lv and Zhang
and monotone mesh patterns with corner shadings,
\newblock
\href{https://arxiv.org/abs/2507.04698}
{arXiv:2507.04698}, 2025.

\end{thebibliography}

\end{document}